% gmpfrxx_mkII Reference Manual
% Compile with: pdflatex gmpfrxx_mkII_manual.tex
%               (or lualatex / xelatex)
%
% Sources:
%   - STATUS.md development log (~8,400 lines, 100+ phases) for policy
%     and history,
%   - public header sources for actual API surface:
%       include/gmpfrxx_mkII/detail/{config,environment,eval_context,
%       expr,integer_conversion,math_mpf,math_mpfc,mpc_environment,
%       mpc_impl,mpf_impl,mpfc_impl,type_traits,zq_impl}.hpp
%   - version_hpp.in.
%   The MPFR class file (mpfr_impl.hpp) was not available; statements
%   about mpfrxx::mpfr_class follow STATUS.md and are flagged where
%   relevant.

\documentclass[a4paper,11pt]{article}

\usepackage[T1]{fontenc}
\usepackage{lmodern}
\usepackage{geometry}
\geometry{margin=25mm}

\usepackage{xcolor}
\usepackage{listings}
\usepackage{hyperref}
\hypersetup{
  colorlinks=true,
  linkcolor=blue!50!black,
  urlcolor=blue!60!black,
  citecolor=blue!50!black,
  pdftitle={gmpfrxx_mkII Reference Manual},
  pdfauthor={Maho Nakata},
}
\usepackage{longtable}
\usepackage{array}
\usepackage{booktabs}
\usepackage{enumitem}
\setlist{itemsep=2pt,topsep=4pt}

% Code listing style.
\definecolor{codebg}{rgb}{0.97,0.97,0.97}
\definecolor{codekw}{rgb}{0.10,0.30,0.70}
\definecolor{codecm}{rgb}{0.40,0.40,0.40}
\definecolor{codestr}{rgb}{0.55,0.10,0.10}

\lstset{
  basicstyle=\ttfamily\small,
  backgroundcolor=\color{codebg},
  keywordstyle=\color{codekw}\bfseries,
  commentstyle=\color{codecm}\itshape,
  stringstyle=\color{codestr},
  language=C++,
  showstringspaces=false,
  breaklines=true,
  columns=fullflexible,
  frame=single,
  framesep=4pt,
  framerule=0pt,
  xleftmargin=8pt,
  xrightmargin=4pt,
  aboveskip=8pt,
  belowskip=8pt,
  morekeywords={mpz_class,mpq_class,mpf_class,mpfc_class,mpfr_class,mpc_class,
                gmp_randclass,mp_bitcnt_t,mpfr_prec_t,mpfr_rnd_t,mpc_rnd_t,
                gmp_randalg_t,mpfr_exp_t,gmpxx,mpfrxx,literals,detail,
                with_precision,real_precision,imag_precision,precision,
                set,set_str,get_str,to_string,to_double,
                real_to_double,imag_to_double,nullptr,constexpr,noexcept,
                override,final,decltype,auto,using,namespace}
}

\newcommand{\code}[1]{\texttt{#1}}
\newcommand{\hdr}[1]{\texttt{#1}}
\newcommand{\sym}[1]{\texttt{#1}}
\newcommand{\envv}[1]{\texttt{#1}}

% Relax line-break tolerance: this manual contains many long monospace
% identifier lists with no natural break points. Without this, LaTeX
% emits overfull \hbox warnings even though the result is acceptable.
\tolerance=2000
\emergencystretch=3em
\hbadness=10000

\title{\textbf{gmpfrxx\_mkII Reference Manual}\\[2pt]
  \large A header-only C++17/20 wrapper for GMP, MPFR, and MPC}
\author{Maho Nakata\\
  \small RIKEN Cluster for Pioneering Research\\
  \small Chiba Institute of Technology}
\date{Version 0.0.1}

\begin{document}
\maketitle

\begin{abstract}
\noindent
\textbf{gmpfrxx\_mkII} is a from-scratch, header-only C++17/20 wrapper for
GMP, MPFR, and MPC. It does not depend on \code{gmpxx.h} or
\code{libgmpxx}, and does not bridge to any pre-existing C++ wrapper.
This manual documents the public API surface as exposed by the project's
detail headers. Internal evaluation machinery (expression-node
evaluators, allocator-aware materialization paths) is mentioned only
where its observable behavior matters to the caller.
\end{abstract}

\tableofcontents
\newpage


% ============================================================
\section{Overview}
% ============================================================

gmpfrxx\_mkII exposes the following arbitrary-precision numeric types:

\begin{itemize}
\item \code{gmpxx::mpz\_class}\,---\,exact integer (GMP \code{mpz\_t}).
\item \code{gmpxx::mpq\_class}\,---\,exact rational (GMP \code{mpq\_t}).
\item \code{gmpxx::mpf\_class}\,---\,GMP MPF floating point.
\item \code{gmpxx::mpfc\_class}\,---\,GMP-only complex, backed by two
  \code{mpf\_class} components. Independent of MPC.
\item \code{mpfrxx::mpfr\_class}\,---\,MPFR correctly-rounded floating point.
\item \code{mpfrxx::mpc\_class}\,---\,MPC correctly-rounded complex.
\item \code{mpfrxx::mpz\_class}, \code{mpfrxx::mpq\_class}\,---\,aliases
  for the corresponding \code{gmpxx} exact types (same type, not subclasses).
\end{itemize}

\subsection{Design properties}

\begin{itemize}
\item Header-only. All public APIs live under \code{include/}.
\item Two public namespaces: \code{gmpxx} (GMP family) and \code{mpfrxx}
  (MPFR/MPC family). Internal facilities live in
  \code{gmpfrxx\_mkII::detail} and are not part of the public API.
\item Strict header separation. Code that needs only GMP can include
  \code{<gmpxx\_mkII.h>} alone and never sees MPFR or MPC headers or
  symbols.
\item Arithmetic operators (\code{+}, \code{-}, \code{*}, \code{/}, unary
  \code{+}, unary \code{-}) return expression nodes, not eager values.
  Materialization happens at construction or assignment.
\item Allocation-aware evaluation: a left-associated addition chain such as
  \code{dst = a + b + c + d} performs zero allocations; a two-sided expression
  such as \code{dst = (a + b) * (c + d)} uses exactly one temporary.
\item Aliasing-safe assignment: when the destination appears in the source
  expression, evaluation falls back to a temporary path.
\item No eager arithmetic operator returning the underlying numeric type is
  defined; the only public surface is the expression-template surface.
\end{itemize}

\subsection{Cross-family arithmetic policy}

\begin{itemize}
\item \code{mpz}, \code{mpq}, and \code{mpf} freely promote within the GMP family.
\item \code{mpz}, \code{mpq}, \code{mpfr}, and \code{mpc} freely promote within
  the MPFR/MPC family. \code{mpfr} promotes to \code{mpc} when at least one
  operand is \code{mpc}.
\item \code{mpf} \emph{must not} be combined with \code{mpfr} or \code{mpc};
  this is enforced at compile time. \code{mpfc\_class} (GMP-only complex)
  must not be combined with \code{mpfr\_class} or \code{mpc\_class} either.
\end{itemize}


% ============================================================
\section{Build and Installation}
% ============================================================

The build system is CMake. The project requires GMP, MPFR, and MPC at
configure time for the full repository; MPFR-only and GMP-only consumers
exist at the include and link-target boundary (see Section~\ref{sec:headers}).

\subsection{Configuration}

\begin{lstlisting}[language=bash]
cmake -S . -B build -DCMAKE_BUILD_TYPE=Release
cmake --build build -j
ctest --test-dir build --output-on-failure
\end{lstlisting}

CMake finder modules \code{cmake/FindGMP.cmake}, \code{cmake/FindMPFR.cmake},
and \code{cmake/FindMPC.cmake} are bundled.

\subsection{Source distribution target}

A \code{dist} target produces a clean source tarball using
\code{git archive HEAD}:

\begin{lstlisting}[language=bash]
cmake --build build --target dist
# produces build/gmpfrxx_mkII-0.0.1.tar.xz
\end{lstlisting}

The archive top directory is \code{gmpfrxx\_mkII-0.0.1/}. The target is
available only when both Git and \code{xz} are detected by CMake. Because
the implementation is \code{git archive HEAD}, uncommitted changes are not
included until committed.

\subsection{Compile-time options}

\begin{description}[leftmargin=*,style=nextline]
\item[\code{GMPFRXX\_MKII\_ASSUME\_FIXED\_PRECISION\_FASTPATH}]
  Define this preprocessor macro before including any wrapper header (or
  via \code{-D}) to enable a fast-path that assumes all
  \code{mpf\_class}, \code{mpfr\_class}, and \code{mpc\_class} operands
  in expressions share the same precision. The macro flips
  \code{gmpfrxx\_mkII::detail::build\_options::assume\_fixed\_precision\_fastpath}
  (a \code{constexpr bool} flag) to \code{true}, which feeds into
  evaluation decisions. Default: off.
\end{description}

\subsection{Version information}

The build system generates a private header \code{detail/version.hpp}
from \code{detail/version.hpp.in}, embedding the Git commit hash. The
following accessors are provided in both \code{gmpxx} and \code{mpfrxx}
namespaces:

\begin{lstlisting}
namespace gmpxx /* and namespace mpfrxx */ {
  constexpr const char* git_commit_hash() noexcept;
  std::ostream& print_git_commit_hash(std::ostream& out);
}
\end{lstlisting}

\noindent
When the version header is unavailable (e.g.\ in a sandboxed build),
\code{git\_commit\_hash()} returns the literal \code{"unknown"}.


% ============================================================
\section{Headers and Link Targets}
\label{sec:headers}
% ============================================================

The library exposes three concentric public headers and four CMake
\code{INTERFACE} targets. Choose the smallest pair sufficient for your
translation unit.

\begin{center}
\begin{tabular}{@{}lll@{}}
\toprule
\textbf{Header} & \textbf{CMake target} & \textbf{Links}\\
\midrule
\code{<gmpxx\_mkII.h>}    & \code{gmpxx\_mkII}    & GMP\\
\code{<mpfrxx\_mkII.h>}   & \code{mpfrxx\_mkII}   & GMP, MPFR\\
\code{<mpcxx\_mkII.h>}    & \code{mpcxx\_mkII}    & MPC (requires MPFR header first)\\
\code{<gmpfrxx\_mkII.h>}  & \code{gmpfrxx\_mkII}  & GMP, MPFR, MPC (aggregator)\\
\bottomrule
\end{tabular}
\end{center}

\subsection{Header rules}

\begin{itemize}
\item \code{<gmpxx\_mkII.h>} is GMP-only. It must not include or expose any
  MPFR or MPC header, type, or symbol. Compile-fail tests enforce this.
\item \code{<mpfrxx\_mkII.h>} provides MPFR (\code{mpfr\_class}) but does not
  expose \code{mpc\_class} or \code{mpfc\_class}.
\item \code{<mpcxx\_mkII.h>} provides \code{mpc\_class} and \emph{requires}
  that \code{<mpfrxx\_mkII.h>} has been included beforehand. It does not
  include MPFR implicitly. A static assertion enforces include order.
\item \code{<gmpfrxx\_mkII.h>} is the convenience aggregator. It pulls in all
  of the above, including \code{<mpcxx\_mkII.h>}.
\end{itemize}

\subsection{Linkage examples}

\begin{lstlisting}[language=bash]
# GMP-only consumer
target_link_libraries(myapp PRIVATE gmpxx_mkII)

# MPFR consumer (no MPC)
target_link_libraries(myapp PRIVATE mpfrxx_mkII)

# MPC consumer (must include both MPFR and MPC headers)
target_link_libraries(myapp PRIVATE mpcxx_mkII)

# Everything
target_link_libraries(myapp PRIVATE gmpfrxx_mkII)
\end{lstlisting}


% ============================================================
\section{Construction and Precision}
% ============================================================

Numeric precision is selected through one of three mechanisms:

\begin{itemize}
\item Implicit, using the family's process-global wrapper default
  (Section~\ref{sec:defaults}).
\item Explicit, by passing \code{mp\_bitcnt\_t} or \code{mpfr\_prec\_t} as
  the second constructor argument.
\item Through the \code{with\_precision(\dots)} static factories on
  \code{mpf\_class}, \code{mpfc\_class}, and \code{mpc\_class} (and
  presumably \code{mpfr\_class}).
\end{itemize}

\noindent
All numeric types reject \code{bool} construction (\code{= delete}).
\code{mpz\_class}, \code{mpq\_class}, and \code{mpf\_class} additionally
accept \code{\_\_int128\_t}, \code{\_\_uint128\_t} where the GMP API allows
exact conversion. Note that \code{\_\_int128} as an \emph{expression-template
scalar leaf} is rejected at compile time; the two paths are distinct.

\subsection{Default construction}

Default construction yields a zero-valued object at the current
wrapper-owned default precision.

\begin{lstlisting}
gmpxx::mpz_class z;        // value 0 (precision is irrelevant for mpz)
gmpxx::mpq_class q;        // value 0/1
gmpxx::mpf_class f;        // 0.0 at default MPF precision
gmpxx::mpfc_class fc;      // (0,0) at default MPF precision per component
mpfrxx::mpfr_class r;      // 0.0 at default MPFR precision
mpfrxx::mpc_class  c;      // (0,0) at default MPC precision per component
\end{lstlisting}

\subsection{Numeric-value constructors}

\subsubsection*{\texttt{mpz\_class}}

\begin{lstlisting}
gmpxx::mpz_class a(42);                 // any non-bool integral type
gmpxx::mpz_class b(__int128_t{...});    // 128-bit signed
gmpxx::mpz_class c(__uint128_t{...});   // 128-bit unsigned
gmpxx::mpz_class d(double{3.0});        // explicit; truncates fractional part
gmpxx::mpz_class e("ff", 16);           // string + base, base default 10
gmpxx::mpz_class f(std::string("123"));
gmpxx::mpz_class g(other_mpz);          // copy
gmpxx::mpz_class h(some_mpz_srcptr);    // explicit, from raw mpz_t
explicit mpz_class(const mpf_class&);   // truncates
explicit mpz_class(const mpq_class&);   // truncates
\end{lstlisting}

\subsubsection*{\texttt{mpq\_class}}

\begin{lstlisting}
gmpxx::mpq_class a(42);                       // integral
gmpxx::mpq_class b(some_mpz);                 // integer numerator
gmpxx::mpq_class c(num_mpz, den_mpz);         // num/den, canonicalized
gmpxx::mpq_class d(355, 113);                 // integral pair
gmpxx::mpq_class e(double{0.5});              // explicit
gmpxx::mpq_class f("355/113");                // string, canonicalized
explicit mpq_class(const mpf_class&);         // exact MPF -> rational
\end{lstlisting}

\subsubsection*{\texttt{mpf\_class}}

\code{mpf\_class} provides numeric-value constructors at the default
precision and overloads taking an explicit precision in
\code{mp\_bitcnt\_t}.

\begin{lstlisting}
gmpxx::mpf_class a(1.5);                      // double, default precision
gmpxx::mpf_class b(1.5, 256);                 // double, 256 bits
gmpxx::mpf_class c(int{42});                  // any non-bool integral
gmpxx::mpf_class d(42L, 256);                 // integral + precision
gmpxx::mpf_class e(some_mpz);                 // exact int -> mpf
gmpxx::mpf_class f(some_mpq, 256);            // exact rational -> mpf at 256
gmpxx::mpf_class g(other_mpf, 256);           // requantize precision
gmpxx::mpf_class h("3.14159");                // string at default precision
gmpxx::mpf_class i("3.14159", 256);           // string + precision
gmpxx::mpf_class j("3.14159", 256, 10);       // string + precision + base
\end{lstlisting}

\paragraph{Constructor diagnostics.} Invalid string input raises
\code{std::invalid\_argument}.

\subsubsection*{\texttt{mpfc\_class} (GMP-only complex)}

\begin{lstlisting}
gmpxx::mpfc_class c(real_mpf, imag_mpf);      // pair of mpf_class
explicit mpfc_class(const mpf_class&);        // real-only
explicit mpfc_class(double{1.5});             // real-only scalar
gmpxx::mpfc_class d(1.5, -1.5);               // scalar pair
gmpxx::mpfc_class e =
    gmpxx::mpfc_class::with_precision(256, 1.0, -1.0);
\end{lstlisting}

\subsubsection*{\texttt{mpc\_class} (MPC complex)}

\code{mpc\_class} has \emph{no} scalar-value constructor. Use
\code{with\_precision} to set the value through doubles, or pass an
\code{mpfr\_class} pair.

\begin{lstlisting}
mpfrxx::mpc_class a(real_mpfr, imag_mpfr);   // pair of mpfr_class
auto b = mpfrxx::mpc_class::with_precision(256);
auto c = mpfrxx::mpc_class::with_precision(256, 256);                // (re,im) prec
auto d = mpfrxx::mpc_class::with_precision(256, /*re=*/1.0, /*im=*/2.0);
auto e = mpfrxx::mpc_class::with_precision(256, 128, 1.0, 2.0);     // mixed prec
\end{lstlisting}

\subsection{The \texttt{with\_precision} factories}

\begin{lstlisting}
// MPF
auto f0 = gmpxx::mpf_class::with_precision(256);
auto f1 = gmpxx::mpf_class::with_precision(256, /*value=*/1.5);

// MPFC: per-component or shared precision
auto c0 = gmpxx::mpfc_class::with_precision(256);
auto c1 = gmpxx::mpfc_class::with_precision(/*re=*/256, /*im=*/128);
auto c2 = gmpxx::mpfc_class::with_precision(256, /*re=*/1.0, /*im=*/-1.0);

// MPC: per-component or shared precision
auto z0 = mpfrxx::mpc_class::with_precision(256);
auto z1 = mpfrxx::mpc_class::with_precision(/*re=*/256, /*im=*/128);
auto z2 = mpfrxx::mpc_class::with_precision(256, /*re=*/1.0, /*im=*/2.0);
\end{lstlisting}

\subsection{Value setters and accessors}

\subsubsection*{Common accessors}

Every numeric class exposes raw GMP/MPFR/MPC pointer access for interop:

\begin{lstlisting}
const mpz_t&  mpz_class::mpz_data() const;     mpz_t&  mpz_data();
const mpq_t&  mpq_class::mpq_data() const;     mpq_t&  mpq_data();
const mpf_t&  mpf_class::mpf_data() const;     mpf_t&  mpf_data();
const mpc_t&  mpc_class::mpc_data() const;     mpc_t&  mpc_data();
// mpfr_class follows the same pattern with mpfr_data().
\end{lstlisting}

\noindent
\code{get\_mpz\_t()} / \code{get\_mpq\_t()} / \code{get\_mpf\_t()}
spellings are also provided for compatibility with the gmpxx idiom.

\subsubsection*{\texttt{mpz\_class} / \texttt{mpq\_class} accessors and setters}

\begin{lstlisting}
mpz_class a;
a.set("42");        a.set(std::string("42"));     // throws on failure
int rc = a.set_str("42", 10);                     // non-throwing; rc==0 on success
std::string s = a.get_str(/*base=*/10);
std::string t = a.to_string(10);
double      x = a.get_d();
unsigned long u = a.get_ui();
signed long   v = a.get_si();
bool ok = a.fits_slong_p();   // also fits_sint_p, fits_sshort_p,
                              //      fits_uint_p, fits_ulong_p, fits_ushort_p

mpq_class q;
q.canonicalize();
mpz_class num = q.get_num();
mpz_class den = q.get_den();
mpz_srcptr np = q.get_num_mpz_t();
mpz_srcptr dp = q.get_den_mpz_t();
double     d  = q.get_d();
\end{lstlisting}

\subsubsection*{\texttt{mpf\_class} accessors and setters}

\begin{lstlisting}
mpf_class f = mpf_class::with_precision(256);
f.set(2.5);
f.set("3.14159");                                  // throws on failure
int rc = f.set_str("3.14159", 10);                 // non-throwing

double        d   = f.to_double();   // also get_d()
unsigned long u   = f.get_ui();
signed long   v   = f.get_si();
mp_bitcnt_t   p   = f.precision();   // also get_prec()

f.set_prec(512);          // change effective precision (preserves value)
f.set_prec_raw(512);      // raw precision change (no value preservation)
f.div_2exp(8);            // f /= 2^8
f.mul_2exp(8);            // f *= 2^8
f.set_epsilon();          // f := 2^-(prec-1)

bool ok = f.fits_slong_p();   // fits_sint_p, fits_sshort_p,
                              // fits_uint_p, fits_ulong_p, fits_ushort_p

std::string s = f.to_string(/*n_digits=*/0);
std::string g = f.get_str();
mp_exp_t exp;
std::string g2 = f.get_str(exp, /*base=*/10, /*n_digits=*/0);
\end{lstlisting}

\subsubsection*{\texttt{mpfc\_class} component access}

\begin{lstlisting}
gmpxx::mpfc_class c(1.0, 2.0);
mpf_class&       re = c.real();    // mutable reference
const mpf_class& ci = c.imag();    // const reference
c.real(other_mpf);                 // setter
c.imag(other_mpf);
mp_bitcnt_t pr = c.real_precision();
mp_bitcnt_t pi = c.imag_precision();
mp_bitcnt_t pp = c.precision();    // == max(pr, pi)
\end{lstlisting}

\subsubsection*{\texttt{mpc\_class} component access}

\begin{lstlisting}
mpfrxx::mpc_class z = mpfrxx::mpc_class::with_precision(256, 1.0, 2.0);
double re = z.real_to_double();
double im = z.imag_to_double();
mpfr_prec_t pr = z.real_precision();
mpfr_prec_t pi = z.imag_precision();
mpfr_prec_t pp = z.precision();    // common precision if equal,
                                   // else min(pr, pi)
mpc_rnd_t rnd = mpfrxx::mpc_class::default_rounding();
\end{lstlisting}

\noindent
For component reads on \code{mpc\_class} use the math-side helpers
\code{mpfrxx::real(z)}, \code{mpfrxx::imag(z)}; both return
\code{mpfr\_class} (Section~\ref{sec:math-mpc}).

\subsection{String construction}

\code{mpz\_class}, \code{mpq\_class}, and \code{mpf\_class} all accept
\code{const char*} and \code{std::string}. Successful \code{mpq\_class}
parses are canonicalized; zero denominators are rejected.

\begin{lstlisting}
gmpxx::mpz_class z("ff", 16);              // 255
gmpxx::mpq_class q("355/113");             // canonicalized
gmpxx::mpf_class f("3.14159", 256);        // explicit precision
mpfrxx::mpfr_class r("1.0e-100", 1024);    // explicit precision
\end{lstlisting}

\paragraph{Throwing vs. non-throwing.} Constructors and assignment
operators that take a string (\code{const char*} or \code{std::string})
throw \code{std::invalid\_argument} on parse failure. The
\code{set\_str()} member returns an \code{int} (0 on success, non-zero on
failure) and preserves the destination value when parsing fails. The
\code{set()} member throws on failure.

\subsection{Increment, decrement, and \texttt{swap}}

\code{mpz\_class}, \code{mpq\_class}, and \code{mpf\_class} support
prefix and postfix \code{operator++} / \code{operator--}. All wrapper
classes provide a non-throwing member \code{swap()} and a free
\code{swap()} that delegates to it.


% ============================================================
\section{Defaults and Environment Variables}
\label{sec:defaults}
% ============================================================

Each numeric family has a process-global, wrapper-owned default. Defaults
are loaded from environment variables at static-initialization time and can
be reloaded or overwritten at runtime. The wrapper does \emph{not} call
\code{mpf\_set\_default\_prec()}; GMP's process-global default is left
untouched.

\subsection{GMP MPF defaults}

\noindent Header: \code{<gmpxx\_mkII.h>}\\
\noindent Initial default precision: \textbf{512 bits}.

\begin{lstlisting}
namespace gmpxx {
  mp_bitcnt_t default_mpf_precision_bits();
  void        set_default_mpf_precision_bits(mp_bitcnt_t bits);
  void        reload_default_mpf_precision_bits_from_environment();
}
\end{lstlisting}

\noindent Environment variable: \envv{MPFXX\_DEFAULT\_PREC\_BITS}.

\subsubsection*{Note on \texttt{mpf\_get\_prec()}}
GMP rounds requested MPF precision up to its internal limb allocation, so
\code{mpf\_get\_prec()} may return a value greater than the requested
default. The wrapper preserves the \emph{requested} default for new
constructions; existing destinations preserve their effective GMP precision
on assignment.

\subsection{MPFR defaults}

\noindent Header: \code{<mpfrxx\_mkII.h>}\\
\noindent Initial default precision: \textbf{512 bits}.\\
\noindent Initial default rounding mode: \code{MPFR\_RNDN}.

\begin{lstlisting}
namespace mpfrxx {
  // Aggregate accessor.
  struct mpfr_default_options {
    mpfr_prec_t precision_bits;
    mpfr_exp_t  emin;
    mpfr_exp_t  emax;
    mpfr_rnd_t  rounding_mode;
  };
  mpfr_default_options default_options();

  // Precision.
  mpfr_prec_t default_precision_bits();
  mpfr_prec_t default_prec();          // alias for default_precision_bits
  void        set_default_precision_bits(mpfr_prec_t);

  // Rounding mode.
  mpfr_rnd_t  default_rounding_mode();
  void        set_default_rounding_mode(mpfr_rnd_t);

  // Exponent range.
  mpfr_exp_t  default_emin();
  mpfr_exp_t  default_emax();
  void        set_default_exponent_range(mpfr_exp_t emin, mpfr_exp_t emax);

  // Reload from environment.
  void        reload_mpfr_defaults_from_environment();
}
\end{lstlisting}

\noindent Environment variables (parsed at startup or on reload):

\begin{center}
\begin{tabular}{@{}ll@{}}
\toprule
\textbf{Variable} & \textbf{Effect}\\
\midrule
\envv{MPFRXX\_DEFAULT\_PRECISION\_BITS} & Wrapper default precision in bits\\
\envv{MPFRXX\_EMIN}                     & Default exponent minimum\\
\envv{MPFRXX\_EMAX}                     & Default exponent maximum\\
\envv{MPFRXX\_ROUNDING\_MODE}           & One of: \code{RNDN}, \code{RNDZ}, \code{RNDU}, \code{RNDD}, \code{RNDA}, \code{RNDF},\\
                                        & \quad or the \code{MPFR\_*} spellings\\
\bottomrule
\end{tabular}
\end{center}

Invalid values fall back to the documented defaults; explicit setters
silently no-op when given out-of-range arguments (precision below
\code{MPFR\_PREC\_MIN}, or \code{emin > emax}). The exponent-range pair
is loaded atomically: if only one of \envv{MPFRXX\_EMIN} or
\envv{MPFRXX\_EMAX} is set, it is applied only if the implied range is
self-consistent. Wrapper expression evaluation guards the MPFR
process-global exponent range with the wrapper-owned defaults via an
RAII guard and restores the previous values afterward.

\subsection{MPC defaults}

\noindent Header: \code{<mpfrxx\_mkII.h>} (the \code{mpc\_class} type itself
lives in \code{<mpcxx\_mkII.h>}).\\
\noindent MPC defaults inherit from MPFR defaults unless
\envv{MPFRXX\_MPC\_*} variables override them. Real and imaginary parts may
be configured independently for both precision and rounding.

\begin{lstlisting}
namespace mpfrxx {
  options_t   default_mpc_options();

  mpfr_prec_t default_mpc_precision_bits();
  mpfr_prec_t default_mpc_real_precision_bits();
  mpfr_prec_t default_mpc_imag_precision_bits();
  void        set_default_mpc_precision_bits(mpfr_prec_t bits);
  void        set_default_mpc_precision_bits(mpfr_prec_t real_bits,
                                             mpfr_prec_t imag_bits);

  mpc_rnd_t   default_mpc_rounding_mode();
  mpfr_rnd_t  default_mpc_real_rounding_mode();
  mpfr_rnd_t  default_mpc_imag_rounding_mode();
  void        set_default_mpc_rounding_mode(mpc_rnd_t rounding);
  void        set_default_mpc_rounding_mode(mpfr_rnd_t real_rnd,
                                            mpfr_rnd_t imag_rnd);

  void        reload_mpc_defaults_from_environment();
}
\end{lstlisting}

\noindent Environment variables:

\begin{center}
\begin{tabular}{@{}ll@{}}
\toprule
\envv{MPFRXX\_MPC\_DEFAULT\_PRECISION\_BITS} & Common default precision\\
\envv{MPFRXX\_MPC\_REAL\_PRECISION\_BITS}    & Real-part precision\\
\envv{MPFRXX\_MPC\_IMAG\_PRECISION\_BITS}    & Imag-part precision\\
\envv{MPFRXX\_MPC\_ROUNDING\_MODE}           & Combined rounding mode\\
\envv{MPFRXX\_MPC\_REAL\_ROUNDING\_MODE}     & Real-part rounding\\
\envv{MPFRXX\_MPC\_IMAG\_ROUNDING\_MODE}     & Imag-part rounding\\
\bottomrule
\end{tabular}
\end{center}


% ============================================================
\section{Expression Templates}
% ============================================================

All public arithmetic operators (\code{+}, \code{-}, \code{*}, \code{/},
unary \code{+}, unary \code{-}) return expression nodes from
\code{gmpfrxx\_mkII::detail} rather than eager numeric values. Materialization
occurs in two contexts:

\begin{enumerate}
\item Construction of a numeric class from the expression.
\item Assignment of the expression to an existing numeric object.
\end{enumerate}

\subsection{Precision policy}

\begin{itemize}
\item \textbf{New materialization} (constructing a fresh object from an
  expression) uses the maximum precision among the relevant object leaves
  in the expression.
\item \textbf{Assignment to an existing object} preserves the destination's
  precision; the expression is evaluated at the destination's precision.
\item Exact \code{mpz}/\code{mpq} leaves contribute \emph{zero} precision to
  mixed exact/MPF or exact/MPFR materialization. The MPF or MPFR leaf
  precision drives the result.
\item For \code{mpc\_class}, real and imaginary precisions are tracked
  separately. New materialization uses the per-component max across leaves;
  assignment preserves per-component destination precision.
\end{itemize}

\subsection{Allocation behavior}

\begin{itemize}
\item Simple leaf-only assignments (\code{dst = a + b}) allocate zero times.
\item Left-associated addition chains (\code{dst = a + b + c + d}) allocate
  zero times.
\item Two-sided compound expressions (\code{dst = (a + b) * (c + d)}) use
  exactly one temporary \code{mpfr\_t} (or \code{mpf\_t}, etc.).
\item Aliasing assignments (where the destination appears in the source)
  fall back to the temporary path; safety is preserved.
\end{itemize}

\subsection{Unary simplification}

A double negation \code{-(-expr)} is simplified to a positive-tag node
\code{pos\_op(expr)} during expression construction. \code{mpz\_class}
additionally supports bitwise complement \code{operator\textasciitilde}.

\subsection{Direct-routed compound assignments}

The \code{mpz\_class} compound assignments \code{+=} and \code{-=} for
expressions of shape \code{mpz * mpz} (or with a small unsigned scalar) are
routed directly to GMP's \code{mpz\_addmul}, \code{mpz\_submul},
\code{mpz\_addmul\_ui}, or \code{mpz\_submul\_ui}. The generic compound
overload is excluded from these shapes to ensure the fast path is used.


% ============================================================
\section{Mixed-Type Promotion}
\label{sec:promotion}
% ============================================================

\subsection{GMP family}

\begin{itemize}
\item \code{mpz + mpz} $\to$ \code{mpz\_class}.
\item \code{mpq + mpq} $\to$ \code{mpq\_class}; \code{mpz + mpq}, division
  in either direction $\to$ \code{mpq\_class}.
\item \code{mpz} or \code{mpq} mixed with \code{mpf} $\to$ \code{mpf\_class}.
  Exact \code{mpz} leaves go through \code{mpf\_set\_z()}. Exact \code{mpq}
  leaves are evaluated by computing numerator and denominator as MPF values
  and dividing (GMP MPF has no \code{mpf\_set\_q()}).
\item \code{mpz} or \code{mpq} mixed with \code{mpfc} $\to$ \code{mpfc\_class}.
\end{itemize}

\subsection{MPFR/MPC family}

\begin{itemize}
\item \code{mpz} or \code{mpq} mixed with \code{mpfr} $\to$ \code{mpfr\_class}.
  \code{mpz} leaves go through \code{mpfr\_set\_z()}; \code{mpq} leaves go
  through \code{mpfr\_set\_q()}.
\item Exact \code{mpz}- or \code{mpq}-result subexpressions embedded in an
  MPFR expression are evaluated in the exact GMP evaluator first, then
  converted once with \code{mpfr\_set\_z()} or \code{mpfr\_set\_q()}. This
  avoids precision loss in intermediate exact arithmetic.
\item \code{mpc} promotes any operand of type \code{mpfr}, \code{mpz}, or
  \code{mpq} when at least one operand is \code{mpc}.
  Conversions go through \code{mpc\_set\_fr()}, \code{mpc\_set\_z()}, and
  \code{mpc\_set\_q()} respectively.
\end{itemize}

\subsection{Forbidden combinations}

The following are rejected at compile time:

\begin{itemize}
\item \code{mpf\_class} combined with \code{mpfr\_class} in either order.
\item \code{mpf\_class} combined with \code{mpc\_class} in either order.
\item \code{mpfc\_class} combined with \code{mpfr\_class} or \code{mpc\_class}.
\end{itemize}


% ============================================================
\section{Scalar Operands}
% ============================================================

Standard built-in scalars participate in expression templates, subject to
strict acceptance rules.

\subsection{Accepted scalars}

\begin{itemize}
\item Signed and unsigned standard integral types, \emph{except}
  \code{bool}.
\item \code{float} and \code{double}.
\end{itemize}

\subsection{Rejected scalars}

\begin{itemize}
\item \code{bool} (compile-fail test enforced).
\item \code{long double} (compile-fail test enforced).
\item \code{\_\_int128} and \code{unsigned \_\_int128} (compile-fail test).
\end{itemize}

\subsection{Internal normalization}

Accepted scalars are normalized at the expression-leaf level to one of three
ABI-stable types:

\begin{center}
\begin{tabular}{@{}ll@{}}
\toprule
\textbf{Source kind}     & \textbf{Normalized leaf type}\\
\midrule
signed integral          & \code{std::int64\_t}\\
unsigned integral        & \code{std::uint64\_t}\\
\code{float}, \code{double} & \code{double}\\
\bottomrule
\end{tabular}
\end{center}

\noindent
\code{std::uint64\_t} max, \code{std::int64\_t} min, and \code{std::int64\_t}
max round-trip exactly through the integer conversion path; the signed
minimum is not negated during conversion.

\subsection{The scalar/scalar boundary}

Expression-template operators always require at least one wrapper object or
wrapper expression operand. Two raw scalars never form an expression node;
this prevents the detail operators from accidentally applying to ordinary
scalar arithmetic.

\subsection{Standard-library interoperability}

\subsubsection*{\texttt{std::common\_type}}

Explicit \code{std::common\_type} specializations are provided for
\code{mpz\_class}, \code{mpq\_class}, \code{mpf\_class}, and
\code{mpfc\_class} paired with each of the following builtin scalar
types, both directions: \code{char}, \code{signed char}, \code{unsigned
char}, \code{wchar\_t}, \code{char16\_t}, \code{char32\_t}, \code{short},
\code{unsigned short}, \code{int}, \code{unsigned int}, \code{long},
\code{unsigned long}, \code{long long}, \code{unsigned long long},
\code{float}, and \code{double}. The result type in each case is the
wrapper class. Cross-class specializations (\code{mpz}/\code{mpq},
\code{mpz}/\code{mpf}, \code{mpq}/\code{mpf}, and the corresponding
\code{mpfc} pairs) follow the promotion table in
Section~\ref{sec:promotion}.

\code{std::common\_type} is also specialized for the expression-template
node types declared in \code{gmpfrxx\_mkII::detail} so that mixing
\code{object\_leaf}, \code{scalar\_leaf}, \code{unary\_expr}, and
\code{binary\_expr} resolves to the underlying \code{result\_type}.

No \code{std::common\_type} specializations are provided on the MPFR/MPC
side (\code{mpfr\_class}, \code{mpc\_class}).

\subsubsection*{\texttt{std::numeric\_limits}}

The following minimal specializations are provided:

\begin{center}
\begin{tabular}{@{}lccccc@{}}
\toprule
\textbf{Type} & \code{is\_specialized} & \code{is\_signed} & \code{is\_integer} & \code{is\_exact} & \code{is\_bounded} \\
\midrule
\code{gmpxx::mpz\_class} & true & true & true  & true  & false \\
\code{gmpxx::mpq\_class} & true & true & false & true  & false \\
\code{gmpxx::mpf\_class} & true & true & false & false & false \\
\bottomrule
\end{tabular}
\end{center}

\noindent
\code{is\_modulo} is \code{false} for all three. No
\code{numeric\_limits} specialization is provided for
\code{mpfc\_class}, \code{mpfr\_class}, or \code{mpc\_class}; for those
types \code{numeric\_limits<T>::is\_specialized} is \code{false}.


% ============================================================
\section{Compound Assignment}
% ============================================================

\subsection{Floating-point family}

\code{mpf\_class}, \code{mpfc\_class}, \code{mpfr\_class}, and
\code{mpc\_class} provide \code{operator+=}, \code{operator-=},
\code{operator*=}, and \code{operator/=} accepting the same operand
families as the binary operators (object, expression, supported scalar).
Compound assignment \emph{always} preserves the destination's precision
regardless of operand precision.

\begin{lstlisting}
mpfrxx::mpfr_class x = mpfrxx::mpfr_class::with_precision(256, 1.0);
mpfrxx::mpfr_class y = mpfrxx::mpfr_class::with_precision(1024, 2.0);
x += y;                       // x stays at 256 bits, evaluated at 256 bits
x *= 3;                       // scalar
x -= (y * y + 1);             // expression RHS
\end{lstlisting}

\code{mpf\_class} additionally supports power-of-two scaling via
\code{operator<<=} and \code{operator>>=} (shift by \code{unsigned
long}), which delegate to \code{mpf\_mul\_2exp} and
\code{mpf\_div\_2exp}.

\subsection{Integer-specific operators}

\code{mpz\_class} exposes the full GMP integer-arithmetic surface as
public operators.

\subsubsection*{Arithmetic compound assignment}

\code{mpz\_class} supports \code{+=}, \code{-=}, \code{*=}, \code{/=}
with operand types \code{mpz\_class}, supported integral and
floating-point scalars, and expression nodes. The
\code{mpz\_class += mpz * mpz} and \code{mpz\_class -= mpz * mpz}
shapes (and scalar-multiplied variants) are routed directly to GMP's
\code{mpz\_addmul} / \code{mpz\_submul} (or \code{mpz\_addmul\_ui} /
\code{mpz\_submul\_ui}) fast paths; the generic compound overload is
specifically excluded from these shapes to ensure dispatch.

\subsubsection*{Modulo}

\begin{lstlisting}
mpz_class operator%(const mpz_class&, const mpz_class&);
mpz_class operator%(const mpz_class&, double);
mpz_class operator%(double, const mpz_class&);
template <Integral I> mpz_class operator%(const mpz_class&, I);
template <Integral I> mpz_class operator%(I, const mpz_class&);

mpz_class& operator%=(mpz_class&, const mpz_class&);
mpz_class& operator%=(mpz_class&, double);
template <Integral I> mpz_class& operator%=(mpz_class&, I);
\end{lstlisting}

\subsubsection*{Bitwise}

\begin{lstlisting}
mpz_class  operator~(const mpz_class&);            // bitwise complement
mpz_class& operator&=(mpz_class&, /* mpz, scalar, or expression */);
mpz_class& operator|=(mpz_class&, /* mpz, scalar, or expression */);
mpz_class& operator^=(mpz_class&, /* mpz, scalar, or expression */);
\end{lstlisting}

\subsubsection*{Shift}

\begin{lstlisting}
mpz_class& operator<<=(mpz_class&, unsigned long);
mpz_class& operator>>=(mpz_class&, unsigned long);
\end{lstlisting}

\subsubsection*{Increment / decrement}

Pre- and post-increment / decrement are defined for \code{mpz\_class},
\code{mpq\_class}, and \code{mpf\_class}, in both lvalue and rvalue
forms.

\subsection{Free functions for \texttt{mpz\_class}}

\begin{lstlisting}
namespace gmpxx {
  mpz_class abs(const mpz_class&);
  mpz_class abs(/* expression node */);   // SFINAE-restricted

  mpz_class sqrt(const mpz_class&);       // floor of integer sqrt

  mpz_class gcd(const mpz_class&, const mpz_class&);
  mpz_class gcd(/* mixed lhs */, /* mixed rhs */);

  mpz_class lcm(const mpz_class&, const mpz_class&);
  mpz_class lcm(/* mixed lhs */, /* mixed rhs */);

  mpz_class factorial (const mpz_class&);
  mpz_class primorial (const mpz_class&);
  mpz_class fibonacci (const mpz_class&);

  // Static spellings on mpz_class:
  static mpz_class mpz_class::factorial (const mpz_class&);
  static mpz_class mpz_class::primorial (const mpz_class&);
  static mpz_class mpz_class::fibonacci (const mpz_class&);
  // plus generic overloads for any non-mpz operand.
}
\end{lstlisting}


% ============================================================
\section{User-Defined Literals}
% ============================================================

User-defined literals are defined under \code{gmpxx::literals} and
\code{mpfrxx::literals}. Re-export into the parent namespace is
\emph{partial}: only the real-valued suffixes are re-exported. The
imaginary-only complex suffixes \code{\_mpfc\_i} and \code{\_mpc\_i}
must be brought in by an explicit \code{using namespace
gmpxx::literals;} or \code{using namespace mpfrxx::literals;}.

\begin{center}
\renewcommand{\arraystretch}{1.15}
\begin{longtable}{@{}>{\ttfamily}l l c >{\raggedright\arraybackslash}p{6.0cm}@{}}
\toprule
\textnormal{\textbf{Suffix}} & \textbf{Type produced} & \textbf{Re-exported?} & \textbf{Notes}\\
\midrule\endhead
\_mpz     & \code{gmpxx::mpz\_class}   & yes & raw / auto-base; supports large hex literals\\
\_mpq     & \code{gmpxx::mpq\_class}   & yes & canonicalized; auto-base\\
\_mpf     & \code{gmpxx::mpf\_class}   & yes & floating and decimal string forms\\
\_mpfc\_i & \code{gmpxx::mpfc\_class}  & no  & imaginary only; needs explicit \code{using namespace gmpxx::literals}\\
\_mpfr    & \code{mpfrxx::mpfr\_class} & yes & floating and MPFR auto-base string forms\\
\_mpc\_i  & \code{mpfrxx::mpc\_class}  & yes & imaginary only\\
\bottomrule
\end{longtable}
\end{center}

\noindent
``Re-exported'' means the suffix is brought into the parent namespace
(\code{gmpxx} or \code{mpfrxx}) automatically by including the relevant
header.

\begin{lstlisting}
using namespace gmpxx;             // brings _mpz, _mpq, _mpf into scope
using namespace gmpxx::literals;   // additionally brings _mpfc_i
using namespace mpfrxx;            // brings _mpfr and _mpc_i into scope

auto z = 0xdeadbeefdeadbeefdeadbeef_mpz;        // large hex via raw parser
auto q = "355/113"_mpq;
auto f = 3.14_mpf;
auto c = 1.0_mpf + 2.0_mpfc_i;                  // GMP-only complex
auto r = 3.14159265358979323846_mpfr;
auto z2 = 1.0_mpfr + (-3.25_mpc_i);             // MPC complex
\end{lstlisting}

Real-only \code{\_mpfc} and \code{\_mpc} literals are intentionally
absent; complex literals follow the standard library's imaginary-only
suffix design (\code{std::complex\_literals}).


% ============================================================
\section{String and Stream I/O}
% ============================================================

\subsection{Exact types (\texttt{mpz\_class}, \texttt{mpq\_class})}

\begin{lstlisting}
namespace gmpxx {
  // string getters
  std::string mpz_class::get_str(int base = 10) const;
  std::string mpz_class::to_string() const;
  std::string mpq_class::get_str(int base = 10) const;
  std::string mpq_class::to_string() const;

  // setters: set_str() is non-throwing and preserves value on failure;
  //         set() throws std::invalid_argument on failure.
  bool mpz_class::set_str(const char* s, int base);
  bool mpq_class::set_str(const char* s, int base);  // rejects 0 denominator
  void mpz_class::set(const char* s, int base = 10);
  void mpq_class::set(const char* s, int base = 10);
}
\end{lstlisting}

\code{std::ostream} insertion supports decimal/hex/octal base selection,
\code{showbase}, \code{uppercase}, \code{showpos}, width, fill, and
alignment.

\subsection{MPF}

\begin{lstlisting}
gmpxx::mpf_class f("1.0e-100", 1024);
f.set("3.14159");
f.set(std::string("3.14159"));
std::string s1 = f.get_str();
std::string s2 = f.get_str(/*exp out*/, /*base*/, /*digits*/);
bool ok = f.set_str("0xff.0p-4", /*base*/ 0);  // value preserved on failure
\end{lstlisting}

\code{std::ostream} insertion uses \code{gmp\_asprintf} and honors basic
ostream formatting flags (precision, \code{fixed}/\code{scientific},
\code{uppercase}, \code{showpoint}, \code{showpos}, width, fill,
alignment). The active C++ locale's decimal point is honored.

\subsection{MPFR}

\code{mpfr\_class} string construction and string assignment throw
\code{std::invalid\_argument} on parse failure, matching MPF behavior.
\code{set\_str()} is non-throwing and preserves the destination value on
failure. \code{std::ostream} insertion supports the standard ostream
formatting flags.

\subsection{MPFC complex}

\code{mpfc\_class} stream insertion materializes \code{mpfc}
expressions through their components and is locale-aware via the
underlying MPF I/O. Equality and inequality comparison
(\code{operator==}, \code{operator!=}) are defined for
\code{mpfc\_class} pairs and for \code{mpfc\_class} \emph{vs.}
\code{mpf\_class}. No ordering operators (\code{<}, \code{<=}, etc.) are
defined.

\subsection{Stream extraction (\texttt{operator>>})}

Stream extraction is implemented for every wrapper type:

\begin{center}
\begin{tabular}{@{}ll@{}}
\toprule
\textbf{Type} & \textbf{Accepted input} \\
\midrule
\code{gmpxx::mpz\_class}    & integer literal in the stream's current base \\
\code{gmpxx::mpq\_class}    & \code{num/den} or integer (rejects zero den.) \\
\code{gmpxx::mpf\_class}    & floating-point literal \\
\code{gmpxx::mpfc\_class}   & \code{(real,imag)} \\
\code{mpfrxx::mpc\_class}   & \code{(real,imag)} \\
\bottomrule
\end{tabular}
\end{center}

\noindent
On parse failure, the destination value (and, for complex types, the
per-component precision) is preserved and the stream's failbit is set.
Raw-pointer overloads of \code{operator<<} and \code{operator>>} are
also defined for \code{mpz\_srcptr} / \code{mpz\_ptr},
\code{mpq\_srcptr} / \code{mpq\_ptr}, and \code{mpf\_srcptr} /
\code{mpf\_ptr}.


% ============================================================
\section{Math Functions}
% ============================================================

All math functions are eager: they accept object or expression operands and
return a freshly constructed numeric object. They do \emph{not} return
expression nodes. Precision policy:

\begin{itemize}
\item Unary functions preserve the materialized operand precision.
\item Binary and ternary functions use the maximum operand precision.
\item Functions taking only scalar arguments come in two overloads: one with
  an explicit destination precision, one using the wrapper default
  precision (e.g., \code{sqrt\_ui(value, prec)} versus \code{sqrt\_ui(value)}).
\item Pair-returning functions (\code{sin\_cos}, \code{sinh\_cosh},
  \code{lgamma}) return both values at the materialized operand precision.
\end{itemize}

\subsection{MPF (\texttt{gmpxx::mpf\_class})}
\label{sec:math-mpf}

Header: \code{<gmpxx\_mkII.h>}.

\begin{description}[leftmargin=*,style=nextline]
\item[Algebraic / power]
  \code{sqrt}, \code{abs}, \code{pow}.
\item[Exponential / logarithm]
  \code{exp}, \code{exp2}, \code{exp10}, \code{expm1};
  \code{log}, \code{log2}, \code{log10}, \code{log1p}.
\item[Trigonometric]
  \code{sin}, \code{cos}, \code{tan},
  \code{asin}, \code{acos}, \code{atan}, \code{atan2}.
\item[Hyperbolic]
  \code{sinh}, \code{cosh}, \code{tanh},
  \code{asinh}, \code{acosh}, \code{atanh}.
\item[Special]
  \code{gamma}, \code{reciprocal\_gamma}.
\item[Constants]
  \begin{sloppypar}
  \code{const\_pi}, \code{const\_e}, \code{const\_log10}, \code{log\_two},
  \code{inv\_log\_two}, \code{log\_ten}, \code{pi\_over\_two},
  \code{pi\_over\_four}, \code{two\_pi}.
  \end{sloppypar}
\end{description}

\paragraph{Implementation notes.}
All listed unary functions evaluate at the destination precision plus
guard bits, then truncate the result back to the destination. Algorithms
used:
\begin{itemize}
\item \code{const\_pi} / \code{const\_e} / \code{log\_two}: hardcoded
  table lookup at 512, 1024, and 2048 bits when the requested precision
  matches; otherwise Gauss--Legendre (\code{pi}) or theta--AGM
  (\code{log\_two}) at the requested precision.
\item Trigonometric functions: argument reduction modulo $\pi/2$ at an
  internal constant precision, followed by Taylor series at working
  precision.
\item \code{log}, \code{log1p}: hybrid path that switches between
  Taylor on $\log(1+x)$ for small $|x|$ and an artanh-series form for
  larger inputs, with cancellation-aware working-precision selection.
\item \code{exp}, \code{expm1}: range-reduced Taylor series.
\item \code{gamma}: Spouge's algorithm with reflection for non-positive
  arguments. The number of Spouge terms is selected from the target
  precision; the function throws \code{std::domain\_error} at
  non-positive integer poles.
\item \code{sqrt}: Newton iteration via GMP's \code{mpf\_sqrt}.
\end{itemize}

\noindent
For applications that require correctly-rounded (last-bit accurate)
results, use \code{mpfrxx::mpfr\_class} (Section~\ref{sec:math-mpfr}),
whose math wrappers delegate directly to MPFR.

\paragraph{Expression-input overloads.}
Every \code{mpf\_class} math function above has a parallel template
overload accepting any expression node whose \code{result\_type} is
\code{mpf\_class}; the expression is materialized first and the
function is then applied.

\subsection{MPFC (\texttt{gmpxx::mpfc\_class})}

Header: \code{<gmpxx\_mkII.h>}.

\begin{description}[leftmargin=*,style=nextline]
\item[Algebraic]
  \code{sqrt}; \code{abs}, \code{norm}, \code{arg}.
\item[Exponential / logarithm]
  \code{exp}, \code{log}.
\item[Trigonometric]
  \code{sin}, \code{cos}, \code{tan},
  \code{asin}, \code{acos}, \code{atan}.
\item[Hyperbolic]
  \code{sinh}, \code{cosh}, \code{tanh},
  \code{asinh}, \code{acosh}, \code{atanh}.
\item[Component / structure]
  \code{conj}, \code{real}, \code{imag}, \code{polar(radius, angle)}.
\item[Power]
  Multiple \code{pow} overloads:
  \code{pow(mpfc, std::uint64\_t)},
  \code{pow(mpfc, std::int64\_t)},
  generic integral overloads, \code{pow(mpfc, mpf\_class)},
  \code{pow(mpfc, mpfc)}, \code{pow(mpf\_class, mpfc)}, and a scalar/mpfc
  overload guarded by \code{is\_supported\_mpf\_scalar\_v}.
\end{description}

\noindent
Each unary function additionally has an expression-input overload
constrained by \code{is\_expression\_node\_v} with
\code{result\_type == mpfc\_class}, materializing the expression first.

\paragraph{Notes.}
\code{abs} returns \code{mpf\_class} (i.e.\ $\sqrt{a^2 + b^2}$),
\code{norm} returns \code{mpf\_class} (i.e.\ $a^2 + b^2$), \code{arg}
returns \code{mpf\_class}, and \code{real}/\code{imag} return
\code{mpf\_class}. \code{conj} returns \code{mpfc\_class}.

\paragraph{Removed.}
\code{mpfc\_class} \code{gamma} and \code{reciprocal\_gamma} were
intentionally removed because the implementation was double-backed via
\code{std::complex<double>} Lanczos and could not honor MPF precision.
They will be reintroduced only with a genuine high-precision
implementation.

\subsection{MPFR (\texttt{mpfrxx::mpfr\_class})}
\label{sec:math-mpfr}

Header: \code{<mpfrxx\_mkII.h>}. All wrappers delegate to MPFR 4.2.2 C APIs
at the materialized precision and the wrapper-owned (or destination)
rounding mode.

\subsubsection*{Algebraic}
\code{sqrt}, \code{sqr}, \code{sqrt\_ui}, \code{rec\_sqrt}, \code{cbrt},
\code{rootn\_ui}, \code{rootn\_si}, \code{agm}, \code{fma}, \code{fms},
\code{fmma}, \code{fmms}, \code{hypot}, \code{min}, \code{max}, \code{dim}.

\noindent (\code{mpfr\_root} is intentionally not wrapped because MPFR
deprecates it; use \code{rootn\_ui} or \code{rootn\_si}.)

\subsubsection*{Power}
\code{pow}, \code{powr}, \code{pow\_si}, \code{pow\_ui}, \code{pow\_z},
\code{pow\_sj}, \code{pow\_uj}, \code{ui\_pow}, \code{ui\_pow\_ui},
\code{compound\_si}.

\subsubsection*{Logarithm and exponential}
\code{log}, \code{log2}, \code{log10}, \code{log1p}, \code{log2p1},
\code{log10p1}, \code{log\_ui};
\code{exp}, \code{exp2}, \code{exp10}, \code{expm1}, \code{exp2m1},
\code{exp10m1};
\code{eint}, \code{li2}.

\subsubsection*{Trigonometric}
\code{sin}, \code{cos}, \code{tan}, \code{sin\_cos};
\code{asin}, \code{acos}, \code{atan}, \code{atan2};
\code{sec}, \code{csc}, \code{cot};
\code{sinu}, \code{cosu}, \code{tanu}, \code{asinu}, \code{acosu},
\code{atanu}, \code{atan2u};
\code{sinpi}, \code{cospi}, \code{tanpi}, \code{asinpi}, \code{acospi},
\code{atanpi}, \code{atan2pi}.

\subsubsection*{Hyperbolic}
\code{sinh}, \code{cosh}, \code{tanh}, \code{sinh\_cosh};
\code{asinh}, \code{acosh}, \code{atanh};
\code{sech}, \code{csch}, \code{coth}.

\subsubsection*{Constants}
\code{const\_pi}, \code{const\_log2}, \code{const\_euler},
\code{const\_catalan}. Each has both an explicit-precision overload and a
default-precision overload.

\subsubsection*{Sign and next-representable}
\code{abs}, \code{signbit}, \code{sgn}, \code{setsign}, \code{copysign};
\code{nextabove}, \code{nextbelow}, \code{nexttoward}.

\subsubsection*{Predicates}
\code{nan\_p}, \code{inf\_p}, \code{number\_p}, \code{integer\_p},
\code{zero\_p}, \code{regular\_p}.

\subsubsection*{Comparison helpers}
\begin{sloppypar}
\code{cmpabs}, \code{cmpabs\_ui}, \code{reldiff}, \code{eq};
\code{greater\_p}, \code{greaterequal\_p}, \code{less\_p},
\code{lessequal\_p}, \code{lessgreater\_p}, \code{equal\_p},
\code{unordered\_p}.
\end{sloppypar}

\subsubsection*{Special functions}
\code{erf}, \code{erfc}, \code{gamma}, \code{lngamma}, \code{digamma},
\code{zeta}, \code{j0}, \code{j1}, \code{y0}, \code{y1}, \code{ai};\\
\code{lgamma} returns \code{std::pair<mpfr\_class,int>} (value plus sign);\\
\code{zeta\_ui} and \code{fac\_ui} provide default- and explicit-precision overloads;\\
\code{jn}, \code{yn}, \code{gamma\_inc}, \code{beta}.

\subsection{MPC (\texttt{mpfrxx::mpc\_class})}
\label{sec:math-mpc}

Header: \code{<mpcxx\_mkII.h>}. All wrappers delegate to the MPC 1.3.1
C API at the materialized operand precision and at
\code{mpc\_class::default\_rounding()}.

\subsubsection*{Unary, returning \texttt{mpc\_class}}
\code{sqrt}, \code{sqr};
\code{exp}, \code{log}, \code{log10};
\code{sin}, \code{cos}, \code{tan};
\code{sinh}, \code{cosh}, \code{tanh};
\code{asin}, \code{acos}, \code{atan};
\code{asinh}, \code{acosh}, \code{atanh};
\code{conj}, \code{proj}.

\subsubsection*{Unary, returning \texttt{mpfr\_class}}
\code{real}, \code{imag}, \code{abs}, \code{norm}, \code{arg}.

\subsubsection*{Binary and ternary}
\code{pow(lhs, rhs)}, \code{agm(lhs, rhs)};
\code{fma(a, b, c)}.

\subsubsection*{Pair-returning}
\code{sin\_cos(z)} returns \code{std::pair<mpc\_class, mpc\_class>}.

\subsubsection*{Constants}
\begin{lstlisting}
mpc_class rootofunity(unsigned long order, unsigned long index);
mpc_class rootofunity(unsigned long order, unsigned long index,
                      mpfr_prec_t precision);
mpc_class rootofunity(unsigned long order, unsigned long index,
                      mpfr_prec_t real_precision,
                      mpfr_prec_t imag_precision);
\end{lstlisting}

\paragraph{Not provided.}
\begin{sloppypar}
MPC 1.3.1 lacks \code{mpc\_gamma} and \code{mpc\_rgamma}; the wrapper
does not provide complex \code{gamma} for \code{mpc\_class}.
Vector-reduction APIs (\code{mpc\_sum}, \code{mpc\_dot}) are not
wrapped. Specialized \code{mpc\_pow\_*} fast paths
(\code{mpc\_pow\_fr}, \code{mpc\_pow\_d}, \code{mpc\_pow\_si},
\code{mpc\_pow\_ui}, \code{mpc\_pow\_z}) are not separately exposed; the
public \code{pow} materializes operands to \code{mpc\_class} and calls
\code{mpc\_pow}. \code{mpc\_mul\_i}, \code{mpc\_mul\_2ui},
\code{mpc\_div\_2ui}, \code{mpc\_mul\_2si}, and \code{mpc\_div\_2si} are
covered by general arithmetic and scalar expression routes rather than
named wrappers.
\end{sloppypar}


% ============================================================
\section{Random Number API}
% ============================================================

There are two random-number wrapper classes, both named
\code{gmp\_randclass} but living in different namespaces with disjoint
purposes.

\subsection{GMP random class (\texttt{gmpxx::gmp\_randclass})}

Header: \code{<gmpxx\_mkII.h>}. RAII wrapper around \code{gmp\_randstate\_t},
non-copyable and non-movable (state ownership).

\begin{lstlisting}
namespace gmpxx {
  class gmp_randclass {
   public:
    // Constructors
    gmp_randclass();                             // default
    explicit gmp_randclass(decltype(gmp_randinit_default));
    explicit gmp_randclass(decltype(gmp_randinit_mt));
    gmp_randclass(decltype(gmp_randinit_lc_2exp_size), mp_bitcnt_t);
    gmp_randclass(decltype(gmp_randinit_lc_2exp), const mpz_class&,
                  unsigned long, mp_bitcnt_t);
    gmp_randclass(gmp_randalg_t, mp_bitcnt_t);   // obsolete form

    // Seeding
    void seed(unsigned long s);
    void seed(const mpz_class& s);

    // Integer generation
    mpz_class get_z_bits(mp_bitcnt_t n);
    mpz_class get_z_bits(const mpz_class& n);
    mpz_class get_z_range(const mpz_class& upper);

    // Float generation (returns an expression node so the caller's
    // destination precision drives generation)
    random_mpf_expr get_f();
    random_mpf_expr get_f(mp_bitcnt_t prec);
    random_mpf_expr get_f(const mpf_class& prototype);
  };
}
\end{lstlisting}

Negative bit counts, non-positive range limits, and over-large LC table
sizes are rejected with errors.

\subsection{MPFR random class (\texttt{mpfrxx::gmp\_randclass})}

Header: \code{<mpfrxx\_mkII.h>}. As MPFR documents, MPFR random generation
runs on a GMP \code{gmp\_randstate\_t}.

\begin{lstlisting}
namespace mpfrxx {
  class gmp_randclass {
   public:
    // Constructors mirror gmpxx::gmp_randclass (default, MT, LC variants,
    // obsolete gmp_randalg_t form).

    void seed(unsigned long s);
    void seed(const gmpxx::mpz_class& s);

    mpz_class get_z_bits(mp_bitcnt_t n);
    mpz_class get_z_bits(const mpz_class& n);
    mpz_class get_z_range(const mpz_class& upper);

    // [0, 1) via mpfr_urandomb (MPF-style)
    random_mpfr_expr get_fr();
    random_mpfr_expr get_fr(mpfr_prec_t prec);
    random_mpfr_expr get_fr(const mpfr_class& prototype);

    // Distributions
    random_mpfr_expr get_fr_uniform();              // mpfr_urandom
    random_mpfr_expr get_fr_uniform(mpfr_prec_t);
    random_mpfr_expr get_fr_normal();               // mpfr_nrandom
    random_mpfr_expr get_fr_normal(mpfr_prec_t);
    random_mpfr_expr get_fr_exponential();          // mpfr_erandom
    random_mpfr_expr get_fr_exponential(mpfr_prec_t);

    // MPFR C-API name aliases
    random_mpfr_expr get_fr_urandomb();   // == get_fr
    random_mpfr_expr get_fr_urandom();    // == get_fr_uniform
    random_mpfr_expr get_fr_nrandom();    // == get_fr_normal
    random_mpfr_expr get_fr_erandom();    // == get_fr_exponential
  };
}
\end{lstlisting}

The expression nodes \code{random\_mpf\_expr} and \code{random\_mpfr\_expr}
integrate with the standard expression-template materialization so that
construction uses the requested or wrapper-default precision and assignment
preserves the destination precision. Non-zero MPFR ternary return values
from \code{mpfr\_urandom}, \code{mpfr\_nrandom}, and \code{mpfr\_erandom}
are rounding information, not failures.


% ============================================================
\section{Examples}
% ============================================================

The \code{examples/} directory contains the following programs. Each is
registered as a CTest target.

\begin{center}
\begin{longtable}{@{}lp{8.5cm}l@{}}
\toprule
\textbf{Program} & \textbf{Demonstrates} & \textbf{Backend}\\
\midrule\endhead
\code{example01}             & Explicit-precision MPF construction, expression-template addition, stream output. & MPF\\
\code{example02\_mpf}        & Default-precision MPF basics. & MPF\\
\code{example02\_mpfr}       & Default-precision MPFR basics. & MPFR\\
\code{example02\_mpfr\_mpc}  & Combined MPFR/MPC use through the aggregator header. & MPFR+MPC\\
\code{example03\_mpf}        & Newton/Heron iteration for $\sqrt{2}$. & MPF\\
\code{example03\_mpfr}       & Newton iteration parallel to the MPF version. & MPFR\\
\code{example03\_mpfc\_math} & Eager MPFC math helpers. & MPFC\\
\code{example04\_*}          & Port of upstream example04 to MPF and MPFR. & MPF / MPFR\\
\code{example05\_*}          & Port of upstream example05. & MPF / MPFR\\
\code{example06\_mpfc}       & Complex example via MPFC. & MPFC\\
\code{example06\_mpc}        & Complex example via MPC pair constructor. & MPC\\
\code{example07\_mpfc}       & ASCII Mandelbrot using MPFC. & MPFC\\
\code{example07\_mpc}        & ASCII Mandelbrot using MPC. & MPC\\
\code{example08\_*}          & Wilkinson polynomial conditioning. & MPFC / MPC\\
\code{example09\_mpfc}       & Aberth--Ehrlich root-finding (default precision). & MPFC\\
\code{example09\_mpc}        & Aberth--Ehrlich root-finding (default precision). & MPC\\
\code{example10\_*}          & Comment-only license-aligned port. & MPF / MPFR\\
\code{example11\_*}          & Upstream example11 port. & MPF / MPFR\\
\code{example12\_*}          & Upstream example12 port. & MPF / MPFR\\
\code{example13\_*}          & Upstream example13 port. & MPF / MPFR\\
\code{example14\_*}          & Upstream example14 port. & MPF / MPFR\\
\code{example15\_mpf}        & BBP-style hexadecimal-digit extraction. & MPF\\
\code{example15\_mpfr}       & BBP-style hexadecimal-digit extraction. & MPFR\\
\bottomrule
\end{longtable}
\end{center}

A minimal program using the aggregator header looks like:

\begin{lstlisting}
#include <gmpfrxx_mkII.h>
#include <iostream>

int main() {
    using namespace gmpxx::literals;
    using namespace mpfrxx::literals;

    auto a = 3.14_mpfr;
    auto b = 2.71_mpfr;
    mpfrxx::mpfr_class c = a * b + 1;
    std::cout << c << '\n';

    mpfrxx::mpc_class z(a, b);
    z = mpfrxx::sin(z);
    std::cout << z.real_to_double() << ' '
              << z.imag_to_double() << '\n';
    return 0;
}
\end{lstlisting}


% ============================================================
\section{Limitations and Notes}
% ============================================================

The following limitations reflect the current state of the library and are
recorded explicitly so users can avoid surprises.

\begin{itemize}
\item \textbf{MPF transcendentals are not correctly rounded.} The
  \code{gmpxx::mpf\_class} math wrappers compute at the destination
  precision plus internal guard bits using arbitrary-precision
  algorithms (see Section~\ref{sec:math-mpf}), but the results are not
  guaranteed last-bit correct. For correctly-rounded transcendentals,
  use \code{mpfrxx::mpfr\_class}.
\item \textbf{MPFC complex \code{gamma} removed.} The previous
  implementation was double-backed via Lanczos and could not honor MPF
  precision. It will be reintroduced only with a genuine high-precision
  implementation.
\item \textbf{MPC \code{gamma} not provided.} MPC 1.3.1 has no
  \code{mpc\_gamma} or \code{mpc\_rgamma}.
\item \textbf{Exact division by zero} follows the underlying GMP behavior;
  no wrapper-specific diagnostic is raised.
\item \textbf{MPF effective precision} as reported by \code{mpf\_get\_prec()}
  may exceed the requested value because GMP rounds up to its internal limb
  allocation.
\item \textbf{Compile-time-rejected expression-template scalar leaves}:
  \code{bool}, \code{long double}, \code{\_\_int128}, \code{unsigned
  \_\_int128}. Note that \code{mpz\_class} accepts \code{\_\_int128\_t}
  and \code{\_\_uint128\_t} \emph{directly} as constructor arguments;
  the rejection applies only when these types appear inside an arithmetic
  expression (e.g.\ \code{mpf + \_\_int128\_t\{...\}}).
\item \textbf{Forbidden mixed families}: \code{mpf} with \code{mpfr};
  \code{mpf} with \code{mpc}; \code{mpfc} with \code{mpfr}; \code{mpfc}
  with \code{mpc}.
\item \textbf{Standard-library traits are partial.}
  \code{std::numeric\_limits} is specialized only for \code{mpz\_class},
  \code{mpq\_class}, and \code{mpf\_class} (with the minimal flags listed
  in Section on standard-library interoperability). For
  \code{mpfc\_class}, \code{mpfr\_class}, and \code{mpc\_class},
  \code{is\_specialized} is \code{false}. \code{std::common\_type} is
  specialized for the GMP-side classes paired with builtin scalars and
  for cross-class promotions; no specializations exist on the MPFR/MPC
  side.
\item \textbf{No MPC vector reductions.} \code{mpc\_sum} and \code{mpc\_dot}
  are unwrapped; ditto \code{mpfr\_sum} and \code{mpfr\_dot}. These require
  a separate public container/API policy that has not been settled.
\item \textbf{Random statistical quality is not certified.} The bundled
  smoke checks catch gross wiring errors only.
\end{itemize}


% ============================================================
\appendix
\section{Public Symbol Index}
% ============================================================

\subsection*{Namespace \texttt{gmpxx}}

\begin{sloppypar}
Types: \code{mpz\_class}, \code{mpq\_class}, \code{mpf\_class},
\code{mpfc\_class}, \code{gmp\_randclass}, \code{random\_mpf\_expr}.\\
Defaults: \code{default\_mpf\_precision\_bits},
\code{set\_default\_mpf\_precision\_bits},
\code{reload\_default\_mpf\_precision\_bits\_from\_environment}.\\
Version info: \code{git\_commit\_hash}, \code{print\_git\_commit\_hash}.\\
Math (\code{mpf\_class}): see Section~\ref{sec:math-mpf}.\\
Math (\code{mpfc\_class}): see Section on MPFC math.\\
Number-theoretic (\code{mpz\_class}): \code{abs}, \code{sqrt},
\code{gcd}, \code{lcm}, \code{factorial}, \code{primorial},
\code{fibonacci}.\\
Literals (in \code{gmpxx::literals}; \code{\_mpz}, \code{\_mpq},
\code{\_mpf} also visible directly in namespace \code{gmpxx},
\code{\_mpfc\_i} only via \code{using namespace gmpxx::literals}):
\code{\_mpz}, \code{\_mpq}, \code{\_mpf}, \code{\_mpfc\_i}.
\end{sloppypar}

\subsection*{Namespace \texttt{mpfrxx}}

\begin{sloppypar}
Types: \code{mpfr\_class}, \code{mpc\_class}, \code{mpz\_class} (alias of
\code{gmpxx::mpz\_class}), \code{mpq\_class} (alias of
\code{gmpxx::mpq\_class}), \code{gmp\_randclass}, \code{random\_mpfr\_expr}.\\
MPFR defaults: \code{default\_options}, \code{default\_precision\_bits},
\code{default\_prec}, \code{set\_default\_precision\_bits},
\code{default\_rounding\_mode},
\code{set\_default\_rounding\_mode}, \code{default\_emin},
\code{default\_emax}, \code{set\_default\_exponent\_range},
\code{reload\_mpfr\_defaults\_from\_environment}.\\
MPC defaults: \code{default\_mpc\_options},
\code{default\_mpc\_precision\_bits},
\code{default\_mpc\_real\_precision\_bits},
\code{default\_mpc\_imag\_precision\_bits},
\code{set\_default\_mpc\_precision\_bits},
\code{default\_mpc\_rounding\_mode},
\code{default\_mpc\_real\_rounding\_mode},
\code{default\_mpc\_imag\_rounding\_mode},
\code{set\_default\_mpc\_rounding\_mode},
\code{reload\_mpc\_defaults\_from\_environment}.\\
Version info: \code{git\_commit\_hash}, \code{print\_git\_commit\_hash}.\\
Math (\code{mpfr\_class}): see Section~\ref{sec:math-mpfr}.\\
Math (\code{mpc\_class}): see Section~\ref{sec:math-mpc}.\\
Literals (in \code{mpfrxx::literals}; \code{\_mpfr} and
\code{\_mpc\_i} also visible directly in namespace \code{mpfrxx}):
\code{\_mpfr}, \code{\_mpc\_i}.
\end{sloppypar}

\subsection*{Macros}

\code{GMPFRXX\_MKII\_ASSUME\_FIXED\_PRECISION\_FASTPATH}\,---\,enable
the fixed-precision evaluation fast path.\\
\code{GMPFRXX\_MKII\_GIT\_COMMIT\_HASH}\,---\,string literal of the
build's Git commit hash, also exposed as
\code{gmpfrxx\_mkII::detail::git\_commit\_hash}.

\subsection*{Environment variables}

\begin{sloppypar}
\code{MPFXX\_DEFAULT\_PREC\_BITS},
\code{MPFRXX\_DEFAULT\_PRECISION\_BITS},
\code{MPFRXX\_EMIN}, \code{MPFRXX\_EMAX},
\code{MPFRXX\_ROUNDING\_MODE},
\code{MPFRXX\_MPC\_DEFAULT\_PRECISION\_BITS},
\code{MPFRXX\_MPC\_REAL\_PRECISION\_BITS},
\code{MPFRXX\_MPC\_IMAG\_PRECISION\_BITS},
\code{MPFRXX\_MPC\_ROUNDING\_MODE},
\code{MPFRXX\_MPC\_REAL\_ROUNDING\_MODE},
\code{MPFRXX\_MPC\_IMAG\_ROUNDING\_MODE}.
\end{sloppypar}

\end{document}
